\section{From the integral forms to the ambient germs}
\label{topo:section}

We now turn the integral isometry of Proposition~\ref{global:integral}
into an equivalence of hypersurface germs.  The two extra prime-power
summands introduced there also ensure that the links are
homotopy spheres.  This permits a direct application of the spherical
knot classification in Theorem~\ref{bg:classification}.

\subsection{The final polynomials and their Seifert forms}

For $j\in\{1,2\}$, put
\begin{equation}
\label{topo:polynomials}
 f_j=q_j^{[1]}+q_j^{[2]}+q_j^{[3]}+q_j^{[4]}
       +z_1^3+\cdots+z_6^3+x^{11}+y^{13}.
\end{equation}
Here $q_j^{[a]}$ denotes a copy of $q_j$ in its own set of variables;
all the displayed summands use disjoint variable sets.  Since $q_j$
has $s_0+2$ variables and $s_0$ is even, the common number of variables is
\begin{equation}
\label{topo:dimension}
 N=4(s_0+2)+6+2=4s_0+16\equiv0\pmod 8.
\end{equation}
In particular, $N>3$.  The weights of the summands, each normalized to
weighted degree one, give positive weights for $f_j$.

Write $(L_j^{\mathrm f},T_j^{\mathrm f},S_j^{\mathrm f})$ for the
Milnor lattice, chosen monodromy operator, and integral Seifert form of $f_j$,
with $T_j^{\mathrm f}=(S_j^{\mathrm f})^{-1}(S_j^{\mathrm f})^{\mathsf t}$.
All monodromy orders under consideration are odd.  Choose a common
odd multiple $m_{\mathrm f}$ of them and an integer $r$ with
$2r\equiv1\pmod {m_{\mathrm f}}$.
As $N$ is even, the symmetric form associated with $S_j^{\mathrm f}$ is
\[
 B_j^{\mathrm f}=S_j^{\mathrm f}(T_j^{\mathrm f})^r,
 \qquad
 S_j^{\mathrm f}=B_j^{\mathrm f}(T_j^{\mathrm f})^{-r}.
\]
The integral Thom--Sebastiani formula identifies these operator--form
data with the data $H_j$ of Proposition~\ref{global:integral}, up to the
same overall sign on the two symmetric forms.  Indeed, every factor
in~\eqref{topo:polynomials}, grouped as four copies of $q_j$, the sum
of six cubes, and the sum $x^{11}+y^{13}$, has even ambient complex
dimension.  Its
symmetrization therefore commutes with taking the indicated tensor
products, and the Thom--Sebastiani signs depend only on these common
dimensions.

Consequently Proposition~\ref{global:integral} supplies an integral
isomorphism $\Psi:L_1^{\mathrm f}\longrightarrow L_2^{\mathrm f}$ with
\[
 \Psi T_1^{\mathrm f}=T_2^{\mathrm f}\Psi,
 \qquad
 \Psi^{\mathsf t}B_2^{\mathrm f}\Psi=B_1^{\mathrm f}.
\]
It also preserves the Seifert forms, because
\begin{equation}
\label{topo:seifert}
 \Psi^{\mathsf t}S_2^{\mathrm f}\Psi
 =\Psi^{\mathsf t}B_2^{\mathrm f}(T_2^{\mathrm f})^{-r}\Psi
 =B_1^{\mathrm f}(T_1^{\mathrm f})^{-r}
 =S_1^{\mathrm f}.
\end{equation}
This is a simultaneous integral isometry, which is precisely the
information needed for the topological step.

\begin{lemma}
\label{topo:analytic}
The polynomials $f_1,f_2$ have isolated critical points at the origin,
are reduced both as polynomials and as analytic germs, and satisfy
$\{\ord_0(f_1),\ord_0(f_2)\}=\{2,3\}$.
\end{lemma}

\begin{proof}
The first nonzero-coordinate argument in Lemma~\ref{seed:chain}
shows that each chain has no critical point anywhere in its affine
space except the origin. The gradient equations of a sum in disjoint
variables separate into the gradient equations of its summands.
The pure powers have the same global property. Consequently the only
critical point of each $q_j$, and then of each $f_j$, is the origin.
The orders of $q_1,q_2$ are two and three.  Disjoint copies have no
monomials that could cancel one another, and all added terms have
degree at least three. The claimed orders follow. More precisely, the
$4E$ balancing squares in each copy of $p_1$ are the only quadratic
terms. Four copies give Hessian ranks $16E$ for $f_1$ and $0$ for
$f_2$, so both ranks are even.

If an irreducible nonunit analytic germ $a$ occurred with multiplicity
at least two in $f_j$, we could write $f_j=a^2b$.  Every first derivative
of $f_j$ would then vanish on $V(a)$.  A nonunit hypersurface germ in
$\C^N$ has dimension $N-1>0$, contradicting the isolated critical
point. Thus $f_j$ is reduced as an analytic germ. The same argument
for a repeated nonconstant polynomial factor would force the gradient
to vanish on its positive-dimensional affine zero set. The global
critical-locus statement just proved excludes this as well, so $f_j$
is reduced as a polynomial.
\end{proof}

\subsection{Spherical links and their embedded equivalence}

Let
\[
 K_j=V(f_j)\cap S^{2N-1}_\varepsilon,
 \qquad
 X_j=S^{2N-1}_\varepsilon\setminus K_j,
 \qquad n=N-1,
\]
where $\varepsilon>0$ is sufficiently small for both polynomials.
The binding $K_j$ is a smooth closed manifold of dimension $2n-1$.
Let $F_j$ be a compact Milnor page, so the fiber of $X_j\to S^1$ is
$\operatorname{int}F_j$.  The inclusion of this interior into $F_j$ is a
homotopy equivalence, and both have the homotopy type of a bouquet of
$n$-spheres.  Moreover $K_j$ is $(n-2)$-connected
by Theorem~\ref{bg:milnor} \cite{Milnor1968}.

\begin{lemma}
\label{topo:sphere}
Each $K_j$ is a homotopy sphere, and each embedding
$K_j\subset S^{2n+1}_\varepsilon$ is a simple knot.
\end{lemma}

\begin{proof}
Every eigenvalue of $T_j^{\mathrm f}$ has the form
\[
 \xi\zeta_{11}^{a}\zeta_{13}^{b},
 \qquad 1\leq a\leq10,\quad1\leq b\leq12,
\]
where the order of $\xi$ divides $3^A\cdot5\cdot7$ for some integer
$A\geq1$.  The factors of orders $11$ and $13$ cannot cancel against
each other or against $\xi$.  Thus every eigenvalue has order divisible
by both $11$ and $13$.

Since $T_j^{\mathrm f}$ is integral and has finite order, its
characteristic polynomial is a product of cyclotomic polynomials
$\Phi_d$, with every occurring index $d$ divisible by $143$.
For $d>1$, the elementary identity
\[
 \Phi_d(1)=
 \begin{cases}
 p,&d\text{ is a power of the prime }p,\\
 1,&d\text{ has at least two distinct prime factors}
 \end{cases}
\]
follows by induction from
$d=\prod_{e\mid d,\ e>1}\Phi_e(1)$: the prime-power divisors already
contribute the full prime factorization of $d$ when $d$ is not a
prime power.  Hence
\begin{equation}
\label{topo:determinant}
 \det(1-T_j^{\mathrm f})=1.
\end{equation}
In particular, $1-T_j^{\mathrm f}$ is an automorphism of the integral
middle homology of $F_j$.

The geometric monodromy acts as $(T_j^{\mathrm f})^{-1}$ in our
convention.  The identity
$1-(T_j^{\mathrm f})^{-1}=-(T_j^{\mathrm f})^{-1}(1-T_j^{\mathrm f})$
shows that its difference from the identity is also invertible over
$\Z$.  The Wang exact sequence of the Milnor fibration now gives
\[
 H_i(X_j;\Z)=
 \begin{cases}
 \Z,&i=0,1,\\
 0,&i\geq2.
 \end{cases}
\]
Indeed, the only possible additional groups are, equivalently, the
kernel and cokernel of $1-T_j^{\mathrm f}$ on $H_n(F_j;\Z)$, and both vanish
by~\eqref{topo:determinant}.  Alexander duality~\cite[Corollary~3.45]{Hatcher2002} therefore makes $K_j$
an integral homology $(2n-1)$-sphere.  Its $(n-2)$-connectivity makes
it simply connected.  The Hurewicz and Whitehead Theorems
\cite[Theorem~4.32 and Corollary~4.33]{Hatcher2002} show that
it is a homotopy sphere.  Since $2n-1\geq5$, the high-dimensional
Poincar\'e theorem included in Theorem~\ref{bg:classification}
makes it homeomorphic to $S^{2n-1}$.

Finally, the homotopy exact sequence of
$\operatorname{int}F_j\to X_j\to S^1$, together with the
$(n-1)$-connectivity of $\operatorname{int}F_j$, gives
\[
 \pi_1(X_j)\cong\Z,
 \qquad
 \pi_i(X_j)=0\quad(2\leq i\leq n-1).
\]
This is the simplicity condition for a $(2n-1)$-knot.
\end{proof}

Orient the links and their Seifert surfaces by the Milnor fibration,
using the same positive push-off convention on both sides.
By~\eqref{topo:seifert}, their integral Seifert matrices are congruent.
Since $n\geq3$, Levine's classification of simple spherical knots
\cite[Theorem~3]{Levine1970}, in the form of
Theorem~\ref{bg:classification}, yields an isotopy of the embedded links. Smooth isotopy extension
then gives an orientation-preserving diffeomorphism of pairs
\begin{equation}
\label{topo:link-equivalence}
 \phi:(S^{2N-1}_\varepsilon,K_1)
       \longrightarrow(S^{2N-1}_\varepsilon,K_2).
\end{equation}
Levine's definition allows smooth homotopy spheres that are not
standard smooth spheres.  Thus no standard-smooth-sphere assertion
is needed here.  Nor is an extra framing hypothesis needed: the
classification is for these oriented embedded knots with their
Seifert forms.  The same implication is part of the classification
of high-dimensional simple fibered links
\cite{Durfee1974,Saeki1999}; see also \cite{Neumann2017}.

\subsection{An explicit homeomorphism of ambient germs}

We finish by spelling out the passage from the sphere to an ambient
neighborhood of the origin.  For a positive weight vector
$w=(w_1,\ldots,w_N)$, define
\[
 a_w(t,x)=(t^{w_1}x_1,\ldots,t^{w_N}x_N),\qquad t>0.
\]
For each nonzero $x$, the function
$t\mapsto\|a_w(t,x)\|^2=\sum_i t^{2w_i}|x_i|^2$
is strictly increasing, with limits zero and infinity at the two
ends of $(0,\infty)$.  Each nonzero weighted orbit therefore meets
$S^{2N-1}_\varepsilon$ in exactly one point.  Consequently
\begin{equation}
\label{topo:weighted-cone}
 A_w:C(S^{2N-1}_\varepsilon)\longrightarrow\overline B^{2N}_\varepsilon,
 \qquad [x,t]\longmapsto a_w(t,x),
 \quad 0\leq t\leq1,
\end{equation}
is a continuous bijection, where all points $(x,0)$ are identified
and their image is zero.  The uniform bound
$\|a_w(t,x)\|\leq\varepsilon t^{\min_i w_i}$ proves continuity
at the cone vertex.  The domain is compact and the target Hausdorff, so
$A_w$ is a homeomorphism.

Let $w^{(j)}$ be the weight vector of $f_j$.  Weighted homogeneity
gives
\[
 f_j(a_{w^{(j)}}(t,x))=t f_j(x).
\]
Thus $A_{w^{(j)}}$ identifies the cone on $K_j$ with
$V(f_j)\cap\overline B^{2N}_\varepsilon$.  Extend the map
in~\eqref{topo:link-equivalence} by
\begin{equation}
\label{topo:ambient-map}
 h(0)=0,
 \qquad
 h\bigl(a_{w^{(1)}}(t,x)\bigr)
      =a_{w^{(2)}}(t,\phi(x))
 \quad(0<t\leq1).
\end{equation}
Equivalently, $h=A_{w^{(2)}}\circ C(\phi)\circ A_{w^{(1)}}^{-1}$,
so $h$ is an ambient ball homeomorphism carrying the first
hypersurface to the second.  It preserves the ball boundary;
restricting to the ball interiors gives the required homeomorphism
germ of ambient pairs.

Together with Lemma~\ref{topo:analytic}, this proves
Theorem~\ref{thm:main}: the reduced hypersurface germs in
\eqref{topo:polynomials} are ambiently homeomorphic and have
multiplicities two and three.

\subsection{Consequences for function germs and initial forms}
\label{topo:consequences}

The same pair has unequal Hessian coranks, as the balancing squares
already show.  Its real coefficients also allow us to strengthen the
ambient equivalence to topological right equivalence.  We then use that
stronger equivalence to compare initial forms after adding one common
square in a new variable.

\begin{corollary}[Negative answer to Arnold's corank problem]
\label{topo:arnold-corank}
In the finite ambient dimension $N>3$ divisible by eight
in~\eqref{topo:dimension}, the polynomials $f_1,f_2$
in~\eqref{topo:polynomials} have different Hessian coranks, although
their reduced complex hypersurface germs are ambiently homeomorphic.
Both polynomials have real coefficients, are weighted homogeneous with
positive rational weights, and have isolated critical points at the origin.
\end{corollary}

\begin{proof}
Theorem~\ref{thm:main} supplies the ambient homeomorphism and the stated
analytic properties. The calculation in the proof of
Lemma~\ref{topo:analytic} gives complex Hessian ranks $16E$ for $f_1$
and $0$ for $f_2$. Since $E>0$ by Lemma~\ref{seed:product}, their
coranks satisfy
\[
 \operatorname{corank}_0(f_1)=N-16E<N=\operatorname{corank}_0(f_2).
\]
\end{proof}

\begin{corollary}[Right-topological multiplicity and parity counterexamples]
\label{topo:right-equivalence}
For the same polynomials $f_1,f_2$ and the single finite ambient dimension
$N>3$ divisible by eight in Theorem~\ref{thm:main}, there is a homeomorphism
germ
\[
 \varphi:(\C^N,0)\longrightarrow(\C^N,0)
 \qquad\text{such that}\qquad
 f_1=f_2\circ\varphi.
\]
Nevertheless, their multiplicities are respectively $2$ and $3$. Thus
neither multiplicity nor multiplicity modulo two is invariant under
topological right equivalence, even for reduced real-coefficient weighted
homogeneous holomorphic germs with positive rational weights and isolated
critical points at the origin.
\end{corollary}

\begin{proof}
For holomorphic germs $f,g:(\C^n,0)\to(\C,0)$ with $n\ge2$ and isolated
critical points at the origin, Saeki's topological $V$-equivalence is
precisely an ambient homeomorphism of their zero-set germs
\cite[Definition~3]{Saeki1989}. Saeki's Corollary~2 identifies this with
topological right-left equivalence \cite{Saeki1989}. King's corollary
\cite{King1978}, recalled in \cite[Section~3, p.~28]{Saeki1989}, then
gives right equivalence up to conjugation. In the precise
source-homeomorphism formulation recorded by Sampaio
\cite[proof of Proposition~2.4]{Sampaio2025Milnor}, there is a homeomorphism germ
$\varphi:(\C^n,0)\to(\C^n,0)$ such that
\[
 f=g\circ\varphi
 \quad\text{or}\quad
 f=(c_1\circ g\circ c_n)\circ\varphi,
\]
where $c_k(z_1,\ldots,z_k)=(\overline z_1,\ldots,\overline z_k)$.
Theorem~\ref{thm:main} supplies these hypotheses for $f=f_1$ and $g=f_2$.
Writing $f_2(z)=\sum_\alpha a_\alpha z^\alpha$ with $a_\alpha\in\R$, we have
\[
 (c_1\circ f_2\circ c_N)(z)
   =\overline{f_2(\overline z)}
   =\sum_\alpha\overline{a_\alpha}z^\alpha
   =f_2(z).
\]
Consequently both alternatives give $f_1=f_2\circ\varphi$.
The multiplicities are $2$ and $3$ by Theorem~\ref{thm:main}, and these
integers are unequal both in $\Z$ and modulo two.
\end{proof}

The homeomorphism $\varphi$ furnished by this argument need not coincide
with the constructed ambient homeomorphism $h$ in~\eqref{topo:ambient-map}.
The parity conclusion concerns multiplicity; the unequal Hessian coranks
above have the same parity.

We now pass to germs in $N+1$ variables by adding the same square to
each polynomial.  The product of $\varphi$ with the identity preserves
right equivalence.  Both new germs have order two, but their quadratic
initial forms have different Milnor fibers.
For a polynomial $g$ vanishing at the origin, write $g_*$ for its lowest
nonzero homogeneous component for ordinary total degree, not the weighted
grading used above. We use the global unit fibre $g_*^{-1}(1)$ as the
initial-form Milnor fibre: it has the homotopy type of the local Milnor
fibre of $g_*$ \cite[Proposition~1.6 in arXiv v3]{BudurEtAl2022ContactLoci}.
This convention retains the full polynomial $g_*$, not its square-free
reduction.

\begin{corollary}[Initial-form Milnor-fibre counterexamples]
\label{topo:tangent-milnor}
In the single ambient dimension $N+1$, put
\[
 \widetilde f_j(z,t)=f_j(z)+t^2,\qquad j=1,2.
\]
The germs $\widetilde f_1,\widetilde f_2:(\C^{N+1},0)\to(\C,0)$ are
topologically right equivalent and have isolated critical points at the
origin, but
\[
 H_0\bigl((\widetilde f_1)_*^{-1}(1);\Z\bigr)\cong\Z,
 \qquad
 H_0\bigl((\widetilde f_2)_*^{-1}(1);\Z\bigr)\cong\Z^2.
\]
Thus integral homology of the initial-form Milnor fibre is not invariant
under topological right equivalence, already in degree zero. In particular,
the initial-form Milnor-fibre homotopy conjecture of Budur, Fern\'andez de
Bobadilla, L\^e, and Nguyen
\cite[Conjecture~1.7 in arXiv v3]{BudurEtAl2022ContactLoci} fails.
\end{corollary}

\begin{proof}
Apply the quadratic case of Sampaio's stabilization
\cite[proof of Proposition~2.3]{Sampaio2025Milnor}.
For the homeomorphism germ $\varphi$ in
Corollary~\ref{topo:right-equivalence}, the product homeomorphism
$\widetilde\varphi(z,t)=(\varphi(z),t)$ satisfies
\[
 \widetilde f_1=\widetilde f_2\circ\widetilde\varphi.
\]
It also carries the corresponding zero-set germs to one another. Moreover,
$\nabla\widetilde f_j(z,t)=(\nabla f_j(z),2t)$, so the critical point
remains isolated on each side.

By the proof of Lemma~\ref{topo:analytic}, the quadratic part of $f_1$ is
\[
 Q=\sum_{\alpha=1}^{16E}u_\alpha^2
\]
in the independent balancing-square coordinates, with $E>0$, whereas
$f_2$ has order three. Hence
\[
 (\widetilde f_1)_*=Q+t^2,\qquad
 (\widetilde f_2)_*=t^2.
\]
Their global unit fibres are respectively
\[
 \left\{v\in\C^{16E+1}:
       \sum_{\alpha=1}^{16E+1}v_\alpha^2=1\right\}
       \times\C^{N-16E},
 \qquad
 \C^N\times\{-1,1\}.
\]
The quadric factor is connected. Indeed, writing $v=x+iy$, its equation
becomes $|x|^2-|y|^2=1$ and $x\cdot y=0$. The map
$(u,y)\mapsto\sqrt{1+|y|^2}\,u+iy$, with
$u\in S^{16E}$ and $y\perp u$, identifies this quadric with the tangent
bundle of the connected sphere $S^{16E}$. The first unit fibre is therefore
connected, whereas the second has exactly two components. This gives the
displayed zeroth homology groups and rules out homotopy equivalence.
\end{proof}

The isolatedness in Corollary~\ref{topo:tangent-milnor} concerns the
stabilized germs, not their initial forms. In particular, the nonreduced
initial form $t^2$ is retained: replacing it by $t$ would change its unit
fibre. No reducedness or isolatedness hypothesis on the initial forms is
part of the cited homotopy conjecture.
